\section{切投影背景：模型集、环面参数与符号因子}\label{app:cut_project_background}

本附录仅作为几何语境的背景引用点：切投影（cut-and-project/model set）、环面相位参数化与 Sturmian 编码的标准事实与定理形式见 \cite{BaakeGrimmAperiodicOrder1,MeyerAlgebraicNumbers,QueffelecSubstitution,LindMarcus1995}。本文的证明链不依赖此处任何陈述，仅在需要时将正文中“旋转扫描/窗口读出/稳定语法”的术语与上述文献口径对齐。

\begin{remark}[phason 参数化与“时间”读出（仅作定位）]\label{rem:app-cutproject-phason-time}
在 cut-and-project 口径中，可将相位视为环面 \(\mathbb{T}=\RR^{d}/\mathcal{L}\) 上的一点，并以一参数平移轨道
\[
\theta(t)=\theta_0+t v\pmod{\mathcal{L}}
\]
表示扫描相位的对象层演化；相关背景见 \cite{BaakeGrimmAperiodicOrder1,MeyerAlgebraicNumbers}。
\end{remark}

\subsection{一维旋转与区间窗读出（Sturmian 接口）}
取圆周旋转 \(T(x)=x+\alpha\ (\mathrm{mod}\ 1)\) 与区间窗 \(W=[0,\beta)\)，二元读出 \(s_t(x)=\ind_W(T^t x)\) 提供正文中“旋转+\(\)窗”的符号接口；其组合性质与平衡性等标准结论参见 \cite{LindMarcus1995}。

